\scanpage{154}
\mspage{V.26}
% archive sheet number: 153

Voici un corollaire du Cas II :

\uline{Corollaire} 5.4 Supposons \inserted{$X$ plat sur $S$ et} $X_0$ loc. une intersection complète et purement de dimension $n$. Alors on a
\[
{\unsure{\underline{\Phi}|\underline{\varphi}}}^i_{\unsure{\nu|v}}=0\quad\text{pour}\quad
\deleted{i\notin[n-d+1,n+1]\quad\unsure{i\leq m'}}\inserted{i\notin[n-d+1,n+1]},
\]
où $d$ est la dimension de l'ens $Z$ des pts de $X_0$ en lesquels $f$ est non lisse, \deleted{donc} et
\[
\Phi^i_\nu=0\quad\text{si}\quad i\notin[n-d+1,n+d+1]
\]
En particulier, si $d=0$, on a
\[
\begin{aligned}
\underline{\Phi}^i_\nu&=0\quad\text{si}\quad i\neq n+1\\
\Phi_\nu&=0\quad\text{si}\quad i\neq n+1,
\end{aligned}
\]
donc l'homomorphisme
\[
H^i(X_0,\Lambda_\nu)\longrightarrow H^i(X_{\bar{\unsure{\eta|y}}},\Lambda_\nu,
\]
est un isomorphisme si $i\neq n,n+1$, un monomorphisme pour $i=n$, un épimorphisme pour $i=n+1$; \inserted{de plus} L'action du grp d'inertie \deleted{\unsure{de}} $I$ sur $H^i(X_{\bar{\eta}},\Lambda_\nu)$ \deleted{est triviale} ou sur $H^i(X_{\bar{\eta}},\mathbb{Z}_\ell)$ est triviale si $i\neq n$.

5.5. Pour terminer le présent paragraphe, nous allons donner une \unsure{variante} des résultats précédents pour la cohomologie non commutative. Supposons $f$ loc. $1$-acyclique \inserted{pour $\unsure{\mathbb{L}|L}$} en toutes les générisations de $x$, \inserted{distinctes de $x$} et \unsure{proposons-nous|proposons} de \deleted{donner} trouver des conditions de \unsure{locale|local} $1$-asphéricité en $x$, pour l'ens de nbs premiers \deleted{$\unsure{\ell\in\mathbb{P}-\{p\}}$}.
\diagram{\begin{tikzcd}[ampersand replacement=\&]
V \arrow[r,no head] \arrow[d,phantom,"\cap" description] \& \bar{V}\arrow[dd,no head]\\
X'\arrow[d] \& \\
S \& \unsure{\bar{S}|\bar{\eta}}\arrow[l]
\end{tikzcd}}
Je dis que pour tt \inserted{$\mathbb{L}$-}grp fini $G$, \deleted{premier à $\mathbb{L}$}, on a
\[
H^i(\bar{V},G)\xrightarrow{\sim}H^i(V_{\bar{\eta}},G)\quad\text{si}\quad i\leq1.
\tag*{(5.5.1)}
\]
